% =====================================================================
\chapter{Testing, Results and Evaluation}
\label{ch:evaluation}

\section{Testing Approach}
\label{sec:eval-testing}
The system was verified at four levels before and during the pilot, in
keeping with the iterative methodology of Chapter~\ref{ch:analysis}:

\begin{description}
  \item[Static verification.] The entire server and mobile codebase is
  written in TypeScript and kept free of type errors; the Prisma ORM
  extends this guarantee to the database boundary, since queries are
  type-checked against the 31-model schema at compile time.
  \item[Component verification.] Each pipeline stage was exercised
  against live data as it was added. The filtering thresholds in
  particular were tuned empirically: rejected fixes are logged with
  their rejection reason (Section~\ref{sec:impl-arbitration}), and
  reviewing those logs during early deployment confirmed that the gates
  rejected genuine noise---urban scatter, impossible jumps---without
  suppressing legitimate motion.
  \item[End-to-end verification on the road.] The complete chain---GPS
  fix, cellular transmission, decoding, arbitration, filtering, ETA,
  push, and on-screen rendering---was verified with live test drives on
  the actual routes, observing the passenger application against the
  visible position of the real bus. Device-configuration changes were
  verified from two independent vantage points (Traccar's report view
  and the server's ingest log, Section~\ref{sec:impl-device}).
  \item[Continuous operational verification.] Because every layer logs
  (FR-13), the running pilot itself acted as a seven-week continuous
  test: anomalies surfaced in the audit data rather than in user
  reports, which is also how the idle-reporting problem of
  Section~\ref{sec:eval-efficiency} was found.
\end{description}

\section{Pilot Deployment and Dataset}
\label{sec:eval-dataset}
\system{} was evaluated in a seven-week pilot, from 6~May to
23~June~2026, on the operational university fleet---real buses, real
routes, real riders, zero simulation. The deployment covered 9 routes
and 3 buses, one of which was equipped with the fixed hardware tracker;
38 passenger accounts were registered, and 71 completed trips were
recorded, during which 46{,}548 raw GPS reports were ingested.
Table~\ref{tab:dataset} summarizes the dataset.

\begin{table}[H]
\caption{Pilot deployment summary.}
\label{tab:dataset}
\centering
\begin{tabular}{@{}lr@{}}
\toprule
\textbf{Quantity} & \textbf{Value} \\
\midrule
Deployment span            & 7 weeks (6 May -- 23 June 2026) \\
Routes                     & 9 \\
Buses                      & 3 (1 fitted with hardware tracker) \\
Registered passengers      & 38 \\
Completed trips            & 71 \\
Raw GPS reports ingested   & 46{,}548 \\
In-trip fixes retained     & 864 \\
Mean GPS fixes per trip    & 37.6 \\
\bottomrule
\end{tabular}
\end{table}

One number in Table~\ref{tab:dataset} demands immediate attention: of
46{,}548 raw reports, only 864 were retained as in-trip fixes. The
remaining ${\sim}98\%$ were transmitted while buses were parked. This
single observation---invisible during design, unmistakable in the
operational data---motivated the efficiency optimization analysed in
Section~\ref{sec:eval-efficiency}, and is in our view one of the most
instructive findings of the whole project.

\section{Measurement Methodology}
\label{sec:eval-method}
All metrics are computed from the production database; no separate test
harness or synthetic workload was used. Position-acquisition latency is
the difference between each fix's device-reported timestamp and its
server-receipt timestamp, both recorded at ingestion
(Section~\ref{sec:impl-traccar}). Source usage is taken from the
canonical source selected per trip. Passenger engagement is measured
from the notification delivery and read records. ETA accuracy is
assessed through the prediction log: every predicted arrival instant is
retained so it can be compared against the corresponding recorded stop
arrival. Because the pilot is small, we report results as a feasibility
demonstration rather than as population statistics.

\section{Results}

\subsection{Position-Acquisition Latency}
\label{sec:eval-latency}
Across 46{,}286 fixes with valid timestamp pairs, the latency from
device fix to server receipt was:

\begin{table}[H]
\caption{Position-acquisition latency (device fix to server receipt).}
\label{tab:latency}
\centering
\begin{tabular}{@{}lr@{}}
\toprule
\textbf{Statistic} & \textbf{Latency} \\
\midrule
Median            & 5.9\,s \\
Mean              & 7.4\,s \\
90th percentile   & 9.9\,s \\
95th percentile   & 11.1\,s \\
\bottomrule
\end{tabular}
\end{table}

A further median 1.2\,s elapsed between server receipt and persistence.
The latency distribution is bounded primarily by the 8\,s Traccar
polling interval rather than by network transport: when the WebSocket
push path is active a fix arrives within about a second, and the polling
path fills the remainder of the distribution. For the passenger this
means the map reflects reality within roughly six seconds in the median
case---comfortably inside the ``live'' perception window for a vehicle
that takes minutes between stops (NFR-1 satisfied).

\subsection{Source Resilience}
\label{sec:eval-sources}
The dual-source design was exercised in earnest rather than remaining a
theoretical safeguard. Of the 70 trips with a recorded canonical source,
45 (64\%) were tracked via the fixed hardware tracker and 25 (36\%) via
the smartphone fallback. The fallback share largely reflects the two
buses not yet fitted with trackers---exactly the deployment reality the
fallback was designed for. Coverage was sustained across the whole
fleet throughout the pilot even though only one bus carried hardware.

Two readings follow. First, graceful degradation works: a bus without a
tracker remained trackable, and mid-trip handovers produced no
passenger-visible artefacts. Second, the design remains hardware-first:
the tracker operates with no driver involvement, and as trackers are
fitted fleet-wide the smartphone path reduces to a rarely needed safety
net (NFR-2 satisfied).

\subsection{ETA Estimation Behaviour}
\label{sec:eval-eta}
The geometric estimator produced a per-stop arrival prediction on every
accepted fix, each carrying its \textsc{high}/\textsc{medium}/\textsc{low}
confidence label. Every prediction was logged against its trip and stop,
so mean absolute error against recorded stop arrivals can be measured
cumulatively as operational data grows. Within the pilot itself the
prediction log was validated end-to-end (predictions recorded, arrivals
recorded, joinable per stop); a statistically meaningful MAE requires
more arrivals per stop than seven weeks provided, and is deliberately
reported as accumulating future evidence rather than overstated from
sparse data.

\subsection{Passenger Engagement}
\label{sec:eval-engagement}
Of 45 arrival and service notifications delivered to riders during the
pilot, 36 (80\%) were read. For a passive daily-utility application this
read rate indicates that the alerts were genuinely attended to rather
than dismissed as noise, supporting the design decision to invest in
proximity-based drop-off alerts (FR-9).

\section{The Idle-Reporting Problem and Optimization}
\label{sec:eval-efficiency}
This section analyses the pilot's most consequential finding.

\subsection{The Discovery}
The fixed-interval tracker transmitted a fix every 5 seconds whenever it
was powered---and it is wired to the vehicle battery, so it was always
powered. Buses, however, spend most of their day parked. The result, in
plain numbers: of 46{,}548 raw reports, only 864 fell inside any active
trip. Roughly \textbf{98\% of all transmissions were waste}---consuming
cellular data on the tracker's SIM, drawing on the bus battery, and
occupying server ingestion for positions nobody would ever look at.

Nothing in the design phase predicted this, because each individual
report is cheap; only the operational aggregate exposed the cost. This
is precisely the kind of finding that distinguishes a deployed
evaluation from a simulation.

\subsection{The Two-Sided Fix}
The waste was attacked at both ends of the pipeline:

\begin{itemize}
  \item \textbf{Device side.} The tracker was reconfigured over SMS
  (\texttt{TIMER,5,300\#}, Section~\ref{sec:impl-device}) to an
  asymmetric cadence: every 5\,s while moving, every 300\,s while
  static---a $60\times$ reduction in parked-state transmissions at the
  source, saving cellular data and battery on the vehicle itself.
  \item \textbf{Server side.} The Traccar polling job was made
  trip-aware: full-rate 8\,s polling now applies only to buses in an
  active trip or pre-trip window, while idle buses are polled every
  60\,s---a $7.5\times$ reduction in idle polling, leaving the WebSocket
  push path and on-trip tracking untouched.
\end{itemize}

\begin{table}[H]
\caption{Idle-reporting efficiency optimization.}
\label{tab:efficiency}
\centering
\begin{tabular}{@{}llrrr@{}}
\toprule
\textbf{Layer} & \textbf{Parameter} & \textbf{Before} & \textbf{After}
& \textbf{Reduction} \\
\midrule
Device (GT06S) & Parked report interval & 5\,s & 300\,s & $60\times$ \\
Server         & Idle poll interval     & 8\,s & 60\,s  & $7.5\times$ \\
\bottomrule
\end{tabular}
\end{table}

\subsection{Verification}
The device-side change was verified from two independent vantage
points: the per-position timestamps in Traccar's report view, and the
server's own ingest log. Both showed parked-state fixes spaced by
minutes rather than seconds after the change, while a live test drive
confirmed that motion still produced dense five-second fixes---that is,
the optimization removed only the redundant reports, not the live
tracking (NFR-4 satisfied). After the optimization the server ingests
essentially only the dense fixes of active trips plus a sparse
five-minute parked heartbeat per bus, in place of the previous
continuous five-second stream.

\section{Requirements Coverage}
Table~\ref{tab:coverage} closes the loop between
Chapter~\ref{ch:analysis} and the pilot.

\begin{table}[H]
\caption{Requirements coverage demonstrated in the pilot.}
\label{tab:coverage}
\centering
\small
\begin{tabular}{@{}lp{10.6cm}@{}}
\toprule
\textbf{Req.} & \textbf{Pilot evidence} \\
\midrule
FR-1--FR-3 & 64\%/36\% tracker/fallback split across 70 sourced trips;
mid-trip handovers without artefacts. \\
FR-4 & Rejected fixes logged with reasons; displayed positions stable
throughout. \\
FR-5 & Per-stop ETA with confidence on every accepted fix; all
predictions logged. \\
FR-6 & Median 5.9\,s acquisition latency; push delivery on every
accepted update. \\
FR-7, FR-8 & Pre-trip vehicle state shown before departure; staleness
flagged within 30\,s. \\
FR-9 & 45 notifications delivered; 80\% read. \\
FR-10 & All pilot riders used the bilingual interface; both scripts
rendered. \\
FR-11--FR-12 & 9 routes, 3 buses, and all schedules administered through
the console. \\
FR-13 & Entire evaluation computed from production logs alone. \\
NFR-1--NFR-7 & Addressed respectively in
Sections~\ref{sec:eval-latency}, \ref{sec:eval-sources},
\ref{sec:impl-deploy}, \ref{sec:eval-efficiency},
\ref{sec:design-security}, \ref{sec:design-realtime},
and \ref{sec:impl-passenger}. \\
\bottomrule
\end{tabular}
\end{table}

\section{Limitations}
\label{sec:eval-limitations}
The results must be read with the pilot's boundaries in mind. The scale
is small---seven weeks, 3 buses, 38 riders---and a subset of trips were
operational tests. Only one bus carried the hardware tracker, so the
64/36 source split reflects fleet hardware coverage as much as tracker
reliability. The GT06S does not report a per-fix accuracy value, so
quality filtering is evaluated at the trip level rather than per fix.
Schedule-adherence analysis was limited by incomplete timetable
association during the pilot. The findings therefore demonstrate
feasibility; they do not claim city-scale statistical generality.
